\paragraph{The two weights seen at a maximal ideal.}
At the chosen label put \(r=k_c+j+t+q\) and number the
coordinates from zero. Let \(A\) be the coefficient algebra
generated by the level matrices and the individual row-weight
projectors in the tensor. For a simple factor \(L\) of the
reduction, take decreasing real numbers \(\xi_{\tau,a}\) at
each label \(\tau\) and choose a joint weight \(\mu\) maximizing
\(\sum_{\tau,a}\xi_{\tau,a}\mu_{\tau,a}\) on \(L\).

If
\(\mu_{\tau,a}<\mu_{\tau,b}\) for \(a<b\), the transposition
of rows changes this sum by
\[
        (\xi_{\tau,a}-\xi_{\tau,b})
                     (\mu_{\tau,b}-\mu_{\tau,a})>0 .
\]
The row permutation preserves the weight set, so the
maximizer is dominant. The highest-weight theorem for this
coefficient algebra embeds \(L\) into the reduced dual Weyl
tensor of weight \(\mu\) and preserves the projectors. At one
label the theorem and the weight bounds give
\[
 \sum_a\mu_a=\sum_a\lambda_a,\qquad
 \lambda_{r-1}\le\mu_a\le\lambda_0,\qquad
 U_a=\lambda_a+r-1-a,\quad U_a'=\mu_a+r-1-a .
\]
In particular \(U,U'\) are strictly decreasing, of equal sum,
with entries in \([0,h_0]\), using \(\lambda_0+r-1\le h_0\).
\label{cor:lift-crystalline-depth-1}

Use the central Capelli generating polynomial
\(C_\nu(X)=\prod_{a=0}^{r-1}(X+\nu_a-a)\).
Its polarizations are coefficients of elementary row additions
followed by projection to the changed weights. The common degree
bound \(|\lambda|+|\mu|\) lets one use the same coefficient
operators on the reduced source and image. On a nonzero image
of \(L\) the two Capelli actions therefore have the same scalar
at every coefficient. The scalars are integral, so this gives
\[
       C_\lambda(X+r-1)\equiv C_\mu(X+r-1)\pmod p,
       \qquad \prod_a(X+U_a)=\prod_a(X+U_a')\ \text{in }k[X].
\]
Use the residue field \(k\) of the coefficient DVR here. Equality
of the two products gives the same multiset of residues, including
when several labels have the same residue.
\label{cor:lift-crystalline-depth-2}

Suppose the forms with values in \(L\) survive at the residual
ideal \(\mathfrak m\). At a chosen double coset \(c\), lift a
reduced value to the integral \(\mu\)-lattice. Its stabilizer
\(H_c\) has order invertible in \(O\). Averaging with the
level action \(R_c\) gives the lift
\[
        v_c^{\rm inv}=|H_c|^{-1}\sum_{h\in H_c} R_c(h)v_c .
\]
The reduction of \(v_c^{\rm inv}\) is the required value, since
the reduced value satisfies the stabilizer equations. Extend these
lifts by covariance. Let \(M_\mu\) be the integral weight-\(\mu\) forms and \(\bar v\) the surviving residual eigenform. Commutation of Hecke translations with reduction gives
\[
 \operatorname{Ann}(M_\mu)\subseteq\operatorname{Ann}(\bar v)=\mathfrak m.
\]
The finite torsion-free Hecke algebra argument
above now gives an integral characteristic-zero eigensystem
over \(\mathfrak m\), realized by a weight-\(\mu\) form since the
coordinate maps preserve independence. A residue-field embedding
aligning the two residual systems lifts to a local Galois
automorphism. This gives the congruence on every generator.
\label{cor:lift-crystalline-depth-3}

\paragraph{Automorphic and crystalline realizations at the chosen place.}
Use the coefficient field and the twisted tuple of this
lift at \(w\mid p\). Choose a \(p\)-adic embedding \(j_w\) inducing \(w\) and an isomorphism \(\iota:\overline{\Q}_p\to\C\). On the coefficient field their composite is prescribed:
\[
 \iota\circ j_w=j_{\mathbb C},
\]
where \(j_{\mathbb C}\) is the selected complex embedding. The latter comparison is algebraic.
The finite Hecke operators can be read through it because
their matrices on the finite class model are formed by
the same right translations and transported level
matrices.

At each member of the selected CM half, use
the finite adelic reading extending that \(p\)-adic
embedding. In the adapted basis the level matrices and
their inverses have integral entries. The twist is a
finite character, so adjoining the roots of unity of its
order puts these matrices and character values in the
same ring of integers \(O\). Evaluation at the finite
double classes gives the common integral form model.
The selected tame level makes the stabilizer orders units
in \(O\), including for the smaller weight \(\mu\).
\label{cor:lift-frame-galois-1}

\begin{proposition}[The oriented crystalline frame and the marked string]
\label{prop:lift-frame-crystalline-reading}
In the selected definite model, with the prescribed algebraic
type and spherical local factors over \(p\), the attached
crystalline representation has the rank-\(r\) normalized
Satake readings. For the \(\lambda\) system assume the specified
length-\(t\) divisor of the Hecke polynomial, \(r<2t\), and the
isolated length-\(t\) run in the weight at every label. Its
unique long Speh factor gives the marked consecutive crystalline
string. After the common character twist the \(\mu\)-frame has the following jumps at the two labels, in the given orientation:
\[
 \text{jump at }\tau=h_0-U'_{\tau,\bar a},\qquad\text{jump at }c\tau=U'_{\tau,a}.
\]
The marked string uses this orientation and the Frobenius transport of the proof.
\end{proposition}

\begin{proof}
\proofstep{The Hecke action after transport to the definite group.}
The characteristic-zero coefficient model first gives
a vector-valued form on the original definite space.
The isometry to the standard definite group transports
its rational and level actions. Use the algebraic
coefficient action at infinity to extend the finite
class values. Pair them using the positive invariant
Hermitian form on the coefficient type and take a
nonzero irreducible component of the resulting weight
space, using an invariant orthogonal complement.
The transported finite Hecke sums are the same on these
components.

In the split reading the prime is the
conjugate of the marked row prime. An isometry of the
two split groups identifies the right coset sums using
the dual double-coset label, which sends the list of elementary-divisor exponents to
\[
 (a_1,\ldots,a_r)\longmapsto(-a_r,\ldots,-a_1).
\]
This reads the given
Hecke eigenvalue, under \(\iota\), at the conjugate
prime. At primes of \(E\) above \(p\) the selected
level is spherical and the level character is trivial
on the corresponding hyperspecial group. The operators
indexed by good split primes away from \(p\) include
all coefficients of their Satake polynomials.
\label{cor:lift-frame-galois-2}

\proofstep{Attachment to the normalized Speh factors.}
Use unitary base change
with these local readings, the discrete spectrum
classification for \(\mathrm{GL}_r\), and Galois
attachment for the regular algebraic polarized cuspidal
constituents. Writing \(F=E^+\) for the totally real field of the definite group, the field hypotheses are
\[
 [F:\Q]\ge2,\qquad E/F\text{ unramified at all finite places}.
\] Base change has the
archimedean infinitesimal characters of the selected
algebraic type and the spherical components over \(p\).
Its discrete factors are Speh representations with
conjugate self-dual unitary cuspidal data. The spherical
component of the induced local product forces the
cuspidal local factors to be spherical.

For a cuspidal datum \(\pi\) of rank \(a\) in a Speh factor
of length \(d\), put
\[
 s_j=\frac{r-1}{2}+\frac{d-1}{2}-j,
 \qquad c_j=\frac{a-1}{2}-s_j \qquad (0\le j<d).
\]
In the unitary infinitesimal parameter let \(z\) range over
the exponents of \(\pi\) at a label. Regular algebraicity of
the full rank-\(r\) parameter gives
\[
 s_j-z\in\Z,
\]
because
\(z+(d-1)/2-j\) occurs there. The same regularity makes the
cuspidal exponents distinct at every label. The twist
\(\pi_j=\pi\otimes |\det|^{c_j}\) has exponents \(z+c_j\)
in \((a-1)/2+\Z\), and is polarized by the norm-power
character descending from \(|\cdot|^{2c_j}\). If \(a\) is odd,
multiply the polarizing character on the real subfield by
the quadratic character of the CM extension to give the
required sign at every real place. Its pullback is trivial.

Apply
Barnet-Lamb--Gee--Geraghty--Taylor, Theorem~2.1.1, to this
regular algebraic polarized cuspidal representation
in rank \(a\). With unitary Satake roots \(\alpha_{\pi,h}\) and Hodge labels
\(z_{\pi,\tau,h}\), the reciprocal normalization in rank \(a\)
reads
\[
 \operatorname{HT}_\tau(\rho_{\pi_j})
     =\{(a-1)/2-z_{\pi,\tau,h}-c_j\}_h
     =\{s_j-z_{\pi,\tau,h}\}_h,\qquad
 \operatorname{Sat}(\rho_{\pi_j})=\{q^{s_j}\alpha_{\pi,h}\}_h
\]
at the indicated good split prime for the Satake expression.
The first equality is the Hodge--Tate normalization and the
second multiset uses \((a-1)/2-c_j=s_j\).
The representation is crystalline above \(p\) because
\(\pi_j\) is spherical there. At a good split prime of
norm \(q\), the twist and the rank-\(a\) local Langlands
normalization multiply the unitary Satake roots by
\[
 q^{-c_j+(a-1)/2}=q^{s_j}.
\]

The direct sum over the cuspidal data and \(j\) therefore
has the full rank-\(r\) Satake roots scaled by
\(q^{(r-1)/2}\), with the corresponding full set of
crystalline jumps. Purity in Theorem~2.1.1 gives equal
absolute values for the unramified roots in each
cuspidal factor. If the cuspidal roots in unitary normalization have the common absolute value \(A>0\), then their product has modulus
\[
 A^a=1,
\]
where \(a\) is the rank of this cuspidal factor. This equality uses the unitary central character. Thus \(A=1\).
\cite{labesse-cm,blgg-weight}\label{cor:lift-frame-galois-3}

\proofstep{Identifying the unique long rank-one string.}
For the selected system of weight \(\lambda\) there is
additional information. Its occurrence in the chosen
block of the row supplies, at a good split prime, the
divisor of its Hecke polynomial coming from the
length-\(t\) string. Its center has modulus one in
unitary Satake normalization. To identify the global
factor, each cuspidal Satake multiset in the discrete
decomposition has unit modulus at this prime, by
temperedness of the cohomological conjugate self-dual
cuspidal factors in their unitary normalization. If a Speh constituent has length \(s\) and cuspidal rank \(a\),
its largest root modulus is \(q^{(s-1)/2}\), and its rank
contribution is \(as\). A root of modulus \(q^{(t-1)/2}\)
therefore requires
\[
       s\ge t,\qquad as\le r<2t .
\]
Hence \(a=1\), and two such constituents cannot fit in rank \(r\).

At each
archimedean label the specified length-\(t\) constant
run in the highest weight has strict drops at its two
ends when those ends are present. The shifted
infinitesimal exponents at this label are distinct.
The long Speh constituent gives a consecutive interval
of exponents. It cannot cross either drop, because each
drop leaves an absent intermediate exponent. The other intervals together have at most \(r-s<t\)
entries. They cannot cover the designated consecutive run of
\(t\) entries. The long interval must meet that run and cannot
cross a missing exponent at either boundary, so \(s\le t\).
Combining the size bound for the constituent and the two boundary gaps gives
\[
 t\le s\le t,\qquad s=t.
\]
Thus the long interval is exactly the specified run. Its Galois factors are the \(t\) successive
crystalline twists of a rank-one character. The
remaining Speh factors give \(\rho\).
\label{cor:lift-frame-galois-4}

\proofstep{Filtration jumps for the two CM orientations.}
The common normalization fixes their readings on both
orientations. Write \(\tau\) for a member of the
\(p\)-adic half chosen by the frame, \(c\tau\) for its
conjugate, \(s_\tau\) for the determinant weight shift
there, and
\[
 \bar a=r-1-a
\]
for a reversed coordinate.
Before the common twist, the type on \(\tau\) has highest
weight \(\mu_{\tau,a}+s_\tau\) and the type on \(c\tau\)
has its dual. The algebraic Hecke character with
crystalline jumps \(h_0-r+1+s_\tau\) on \(\tau\) and
\(-s_\tau\) on \(c\tau\) is chosen unramified over \(p\).
The sum of the two jumps is constant on all pairs, as
required for the algebraic character. Under the
reciprocal Frobenius convention the untwisted jump is
\(\bar a-\lambda_{{\rm aut},\bar a}\) in either local
highest weight \(\lambda_{\rm aut}\).

Adding the character
jump gives the two arrays
\[
 J_{\tau,a}=h_0-U'_{\tau,\bar a},\qquad
 J_{c\tau,a}=U'_{\tau,a},\qquad
 U'_{\tau,a}=\mu_{\tau,a}+r-1-a .
\]
Thus the marked orientation for \(U'\) is the conjugate
of the half used for the complex automorphic frame.
Apply the same twist to the selected \(\lambda\) system.
On that orientation the marked line jumps are
\(B_\tau+s\), \(0\le s<t\), with \(B_\tau\) the bottom
of the shifted run of \(\lambda\). On the other orientation
their bottom is \(h_0-B_\tau-(t-1)\), with the order of
the marked indices reversed. Both arrays and both
families of lines therefore lie in \([0,h_0]\).
The Frobenius permutation acts on all embeddings of the
unramified host and commutes with their CM involution.
This determines the transported label and its
orientation together in the crystalline comparison.
\label{cor:lift-frame-galois-5}
\end{proof}


\begin{lemma}[Identifying the marked factors in the companion reduction]
\label{lem:companion-residual-factors}
On the common coefficient extension of the two normalized
Hecke systems, the residual semisimplification is
\[
 \bar r_\mu^{\rm ss}\simeq \bar\rho^{\rm ss}\oplus\bigoplus_{s=0}^{t-1}\bar\chi_s.
\]
One may keep the marked occurrence of each scalar line in a
composition series of its actual reduced lattice.
\end{lemma}

\begin{proof}
At the good degree-one Frobenius classes, congruence of
the two Hecke systems gives
\[
 \det(X-r_\mu(g))\equiv
 \det(X-\rho(g))\prod_{s=0}^{t-1}(X-\chi_s(g))
                 \pmod {\mathfrak m_{\overline{\Q}_p}} .
\]
Indeed the reversed polynomial has Hecke coefficients times
signed integral norm powers and powers of the common twisting
character, which has unit values by compactness. For degree \(r\), reversing by
\[
 P(X)\mapsto X^rP(X^{-1})
\]
is a permutation of the coefficient list. It therefore
preserves the congruence after the common unit normalizations. The congruence
locus is closed because all coefficients are continuous and
the open unit ball in an ultrametric field is closed.
Degree-one Chebotarev density after finitely many omissions
extends it to every group element, also for the reciprocal
Frobenius convention. Choose a common finite coefficient field
and stable lattices.

For the two residual representations \(\bar R_1,\bar R_2\), Brauer--Nesbitt gives
\[
 \det(1-X\bar R_1(g))=\det(1-X\bar R_2(g))\ (\forall g)\ \Longrightarrow\ \bar R_1^{\rm ss}\simeq\bar R_2^{\rm ss}.
\]
A composition series therefore has the
prescribed line characters among its factors with their
multiplicities. Keep their positions in the filtration.
The block comparison takes these lattices over a common
finite coefficient extension with the marked characters
defined there. Compactness gives an invariant lattice
by taking the span of translates of one integral basis.
The coefficient field can be enlarged to contain the
labeled fields and the primitive \(p\)-th roots used in
the cyclotomic frame. This computation uses the original
residual filtration on the \(r_\mu\) side, with its
marked factors identified by the semisimplification.
\label{cor:lift-crystalline-depth-4}
\end{proof}

\paragraph{What the residual lines imply about the head.}
\hypertarget{source-lift-weight-exchange}{}
Let \(H=h_0\) and \(K=\lfloor H/(p-1)\rfloor\).
Here \(p>y>8\), \(p\) is unramified in the CM field, and
the large-prime inequalities give \(t+2K<p\).
Use all embeddings of the CM field at \(p\). At the conjugate embedding use the decreasing sequence obtained by reflection and reversal. In zero-based coordinates the operation is
\[
 (u^c)_a=H-u_{r-1-a}\quad(0\le a<r),\qquad (u^c)^c=u .
\] The
bottom of the new block is \(H-(B+t-1)\). Write
\(\sigma i\) for the source label of the Frobenius matrix
ending at \(i\).

In the local cyclic indexing it is \(i-1\).
This permutation commutes with conjugation. The Galois attachment can use either
orientation for the whole family. Pulling back the two
sequences and the block along that involution identifies its
\(J\) with \(U'\), and its line weights with
\(b_{i,s}=B_i+s\) at each embedding \(i\). A calculation at
the transported label then reads at \(i,\sigma i\) in the
original labels. Reflection preserves the characteristic-polynomial comparison. Indeed, for any list of residues of length \(r\),
\[
 \prod_{a=0}^{r-1}\bigl(X-(H-u_a)\bigr)=(-1)^r\prod_{a=0}^{r-1}\bigl((H-X)-u_a\bigr).
\]

\begin{proposition}[A depth bound from the residual lines]\label{prop:residual-head-depth}
For the shifted highest weights, conjugate labels and consecutive
line block just specified, the deficit at a head cut satisfies
\[
 \sum_{a<k_c+j}(\lambda_{i,a}-\mu_{i,a})
   \le qH+(p m_{\sigma i}+m_i)(t+2K).
\]
Here \(m_i\) counts the original shifted weights outside the consecutive
block at \(i\) whose residues occur in the block enlarged by \(K\)
at each end. The count \(m_{\sigma i}\) uses the Frobenius source label. At the indices satisfying the shape and
prime exclusions of Section~\ref{subsec:lift-exponent-primes},
both corrections vanish and the depth is at most \(\lfloor qh\rfloor\),
uniformly over the occurring Hecke ideal and the weight \(\mu\).
\end{proposition}

\begin{proof}
Use the norm coordinate of Section~\ref{subsec:cut-wach-models},
written \(X=Z\) in this calculation, with the same terminal
Hodge labels. Apply its construction to
the lattice of the congruent weight \(\mu\), to the inverse
images of the residual flag and to the specified scalar lines.
Their coefficient extension is over the selected place. Their
jumps are respectively \(U'_{i,a}\) and \(b_{i,s}\) and lie
in \([0,H]\). Write \(F\) for the residue field. In the resulting
\(F[[X]]\) lattice the reduced Frobenius has an integral right
inverse up to \(X^H\). For any \(d\)-rowed minor \(M\) of the reduced Frobenius matrix the elementary divisors give
\[
 \operatorname{ord}_X M\ \ge\ s_i(d),\qquad s_i(d)=\sum_{a=r-d}^{r-1}U'_{i,a}.
\]
The first cyclotomic matrix coefficient \(\Theta_i\) has
characteristic roots \(-\bar c U'_{i,a}\).
The residual filtration gives the same ranks and marked lines in
\(F((X))\) by exact 'etale realization.

Its intersection with
\(F[[X]]^r\) is a saturated flag, so extending bases at successive
steps gives one basis adapted to it. Each marked line occupies its own
diagonal block. Its two chosen generators in the Laurent
line differ by a unit times \(X^{\kappa_{i,s}}\), with the
sign fixed so that Frobenius in the adapted basis has order
\[
        \alpha_{i,s}=b_{i,s}+p\kappa_{\sigma i,s}-\kappa_{i,s}.
\]
In fact Frobenius sends the parameter to \(X^p\), and the
reduced rank-one Frobenius is a unit times \(X^{b_{i,s}}\).
The latter follows from the rank-one determinant, whose
order before reduction is \(b_{i,s}\) at \(X-p\) by height
and the local elementary divisor. Integrality gives
\(\alpha_{i,s}\ge0\) and the height inverse gives
\(\alpha_{i,s}\le H\) on the singleton block. The inverse
is block triangular after inverting \(X\), so the two
diagonal entries multiply to \(X^H\).

\proofstep{From the determinant to the minor estimate.}
To constrain individual head weights we pass from the determinant to a minor.

For the cyclotomic generator, \(\gamma(X)/X\) has first
terms \(1+\bar c X\) modulo \(X^2\) in this reduced
coordinate. For a unit \(v=v_0+v_1X+O(X^2)\), the constant and linear terms cancel in the quotient:
\[
 \gamma(v)\equiv v_0+v_1X\equiv v\pmod{X^2},\qquad \gamma(v)/v\equiv1\pmod{X^2}.
\] The coefficient of \(X\) on the same marked line
is therefore \(-\bar c(b_{i,s}-\kappa_{i,s})\).
An integral basis change conjugates this first coefficient by
its constant matrix, since the cyclotomic action fixes the
change of basis modulo \(X^2\).

Taking characteristic
polynomials of the block triangular matrix and cancelling
the nonzero common scalar shows that the residues
\(b_{i,s}-\kappa_{i,s}\), counted with multiplicity, occur
among those of \(U'_i\). The determinant of the Frobenius
minor on the marked positions is the product of their
diagonal entries. The minor divisibility then gives
\begin{equation}\label{eq:lift-line-minor}
       s_i(t)\le\sum_s b_{i,s}
              +p\sum_s\kappa_{\sigma i,s}-\sum_s\kappa_{i,s}.
\end{equation}
The flag and the rank-one comparison here come from the
same residual representation at the same maximal ideal.
\label{cor:lift-weight-exchange-1}

Since \(b_{i,s}\) and \(\alpha_{i,s}\) belong to \([0,H]\),
we have \(|p\kappa_{\sigma i,s}-\kappa_{i,s}|\le H\).
For fixed \(s\), let \(j\) maximize \(\lvert\kappa_{j,s}\rvert\).
Apply the relation at \(i=\sigma^{-1}j\), where the term
with coefficient \(p\) is \(\kappa_{j,s}\). The triangle inequality at this predecessor gives
\[
 p|\kappa_{j,s}|\le H+|\kappa_{\sigma^{-1}j,s}|
            \le H+|\kappa_{j,s}| .
\]
Hence every \(|\kappa_{i,s}|\le H/(p-1)\).
To bound the sum of the shifts, enlarge the block
\(I=[B_i,B_i+t-1]\) by \(K\) at both ends to \(I^+\).

Let \(m_i\) count original weights outside \(I\) whose
residue occurs in \(I^+\). The strict length bound makes
reduction modulo \(p\) injective on the integers in \(I^+\).
The \(t\) shifted entries \(b_{i,s}-\kappa_{i,s}\) lie
there. Match their residues to distinct occurrences among
the \(U_i\) using the two multiset comparisons. At most
\(m_i\) are matched outside \(I\). All other matches are
equalities in \(I^+\). Pair the remaining shifted entries
arbitrarily with the unused block entries. Each difference
is bounded by \(t+2K\), so
\[
                \left|\sum_s\kappa_{i,s}\right|
                            \le m_i(t+2K).
\]
This also holds at \(\sigma i\), where the exceptional count
uses its own block. Reflection preserves the count because it
carries the block, its enlargement and residue congruence
bijectively to their reflected versions.

\proofstep{Returning to the head of the weight.}
The deficit of the first \(k_c+j\) weights equals the excess
of the rest by equality of total sums. Of the last \(t+q\)
entries of \(U'_i\), the smallest \(t\) have sum \(s_i(t)\)
and the other \(q\) are at most \(H\) apiece. In \(U_i\)
the corresponding total is at least \(\sum_s b_{i,s}\)
because the bottom \(q\) entries are nonnegative. Equation~\eqref{eq:lift-line-minor}
thus proves
\[
 \sum_{a<k_c+j}(\lambda_{i,a}-\mu_{i,a})
   \le qH+(p m_{\sigma i}+m_i)(t+2K).
\]
The staircase cancels term by term on the left. The shape
and prime exclusions of Section~\ref{subsec:lift-exponent-primes}
give \(m_i=m_{\sigma i}=0\). Every surviving highest
weight at this index with a larger weight therefore has depth at most
\(i_i=\lfloor qh\rfloor\), with the choice of
\(\mu\) and of the occurring Hecke ideal unrestricted.
\label{cor:lift-weight-exchange-2}
\end{proof}

